\subsection{Locally finite coefficients and arbitrary powers of a finite ring}
\label{C:locally-finite-subsection}

Fix the split orthogonal group $O_m$ over $\mathbb Z[1/2]$, and put
$\mathcal W(R)=WG_{\mathrm{sp}}(BO_m/R)$ for a commutative
$\mathbb Z[1/2]$-algebra $R$. A ring is \emph{locally finite} if every
finite subset is contained in a finite unital subring. The finite-base
answer extends to such rings with an exact finite-stage equality
criterion. For powers of one finite ring it gives an explicit
description with no remaining representation-isomorphism quotient.

\begin{theorem}[Proved: filtered continuity with explicit finite-stage maps]
\label{C:splitWG-filtered-colimits}
For every filtered system $(R_i)$ of commutative rings with $2$
invertible, scalar extension induces an isomorphism
\[
 \mathop{\rm colim}_i\mathcal W(R_i)
       \ \xrightarrow{\ \sim\ }\ 
 \mathcal W\bigl(\mathop{\rm colim}_iR_i\bigr) .             \tag{C.3.L1}
\]
It preserves products and every exterior operation. Two classes
represented at finite stages have the same image if and only if
their scalar extensions are equal at some common later stage.

Consequently, for every locally finite $R$ with $2$ invertible,
$\mathcal W(R)$ is the explicitly specified directed limit of the
finite-base groups and operations of
Theorem~\ref{C:Om-finite-base-algorithm}, indexed by the finite unital
subrings of $R$. Every transition map is computed by the finite
matrix procedures described below. This is an exact finite-witness
equality criterion; it does not assert that an arbitrary presentation
of a locally finite ring permits one to decide whether such a later
witness exists.
\end{theorem}
\begin{proof}
The ordinary nonequivariant continuity principle appears, under
Noetherian limit hypotheses, in
\cite[Appendix B, Lemma 49 and Corollary 51, pp.26--27]{C-BachmannContinuity}.
We give the full finite-data argument for the present split equivariant
group and arbitrary filtered systems.

Write $R=\mathop{\rm colim}_iR_i$. A finite projective $R$-module
$P$ is the image of an idempotent matrix $e\in M_N(R)$.
Its finitely many entries lift to one $R_i$; the finitely many
identities in $e^2=e$ hold after increasing $i$.
Thus $P$ descends to a finite projective module.
The coordinate Hopf algebra
\[
 H_R=R[x_{ab}:1\leq a,b\leq m]/(X^tJX-J)
\]
is obtained by scalar extension from $\mathbb Z[1/2]$ and commutes
with filtered colimits, as does $H_R\otimes_RH_R$.
A coaction on $P=eR^N$ is specified by a finite matrix over $H_R$,
supported by the source and target idempotents. Choose polynomial
representatives for all its entries. The support, counit and
coassociativity conditions are finitely many equalities in $H_R$
and $H_R\otimes_RH_R$. All coefficients lift to one stage and all
equalities hold at a common later stage: an element is zero in a
filtered colimit precisely when its image is zero later. This
argument uses no flatness of a transition map and no freeness
assertion for a bounded polynomial-coefficient module.

A nonsingular symmetric form is an isomorphism $q:P\to P^*$.
Relative to the idempotent presentations of $P$ and $P^*$ it is
specified by finite matrices for $q$ and its inverse. Their support
conditions, two inverse identities and the symmetry identity are
finite equations over $R$. Invariance under the coaction is a
finite matrix identity over $H_R$. These data and equations
therefore descend after increasing the stage once more.
The same argument applies to an equivariant isometry and its
inverse, and to an equality between two such maps. In particular,
two descended forms that become isometric over $R$ are isometric
at a common finite stage.

Let $\mathcal M(R)$ denote the commutative monoid of isometry
classes under orthogonal sum. The preceding paragraph proves
$\mathcal M(R)=\mathop{\rm colim}_i\mathcal M(R_i)$.
One can now take group completions. Explicitly, equality
$[Q]-[P]=[Q']-[P']$ in the group completion means that an actual
form $T$ exists with
\[
 Q\perp P'\perp T\simeq Q'\perp P\perp T.
\]
The form $T$ and this isometry both descend. This proves
injectivity, surjectivity and the asserted finite-stage equality
criterion in (C.3.L1). It avoids any cancellation assumption about
the monoid.

Tensor products and their forms commute with scalar extension.
So do $\Lambda^rP$, its coaction and its determinant form: the
exterior construction on a finite projective summand is obtained
by applying $\Lambda^r$ to its idempotent, and the form uses
determinants. No division by $r!$ is involved. On virtual classes
the formula
$\lambda_z([Q]-[P])=\lambda_z(Q)\lambda_z(P)^{-1}$
then proves compatibility of every coefficient.

Finally, the finite unital subrings of a locally finite ring form
a directed system. Their characteristic is odd because $2$ is
invertible in $R$: if the prime subring had even characteristic,
multiplication by $2$ would have a nonzero kernel in $R$.
Thus the inverse of $2$ already lies in the finite prime subring.
All these finite subrings are eligible for
Theorem~\ref{C:Om-finite-base-algorithm}.
For an inclusion $A\subset B$ of two of them, extend the entries
of a chosen coaction and Gram matrix to $B$. Decompose $B$ into
its computable finite local factors; then apply the canonical
coaction, indecomposable and finite-form procedures of that theorem.
This computes the image of every concrete generator, including a
two-torsion generator represented as the difference of two forms.
Extend additively to virtual classes. Products and exterior
operations are already computed at finite stages, completing the
description of the direct system.
\end{proof}

\begin{theorem}[Proved: a quotient-free answer for powers of a finite ring]
\label{C:finite-ring-powers}
Let $A$ be any finite commutative ring with $2$ invertible and let
$I$ be any set. Coordinate evaluation gives an isomorphism
\[
 \mathcal W(A^I)\ \simeq\
 \bigl\{f:I\longrightarrow\mathcal W(A):
                    f(I)\text{ is finite}\bigr\}.           \tag{C.3.L2}
\]
Addition, multiplication and all exterior operations on the right
are pointwise. The values in $\mathcal W(A)$ have the explicit
canonical indices and coordinates of
Theorem~\ref{C:Om-finite-base-algorithm}.

A unique description of an element is therefore its finite list of
distinct nonzero values $b$, together with their disjoint nonempty
level sets $I_b\subset I$; the function is zero on the complement.
This leaves no unclassified representation or form isomorphism
quotient. On inputs given by finite partitions and their values,
the operations are finite procedures using Boolean operations
and equality for the represented subsets of $I$, and the
finite-base procedures on the values. For an arbitrary set $I$,
an effective representation of its subsets is part of this
computational access model, not an implicit assumption.
\end{theorem}
\begin{proof}
For a finite partition $\Pi$ of $I$ into nonempty cells, let
$A^\Pi\subset A^I$ be the subring of functions constant on each
cell. It is a finite product of copies of $A$. The partitions
are directed by common refinement, and
\[
 A^I=\mathop{\rm colim}_{\Pi}A^\Pi.
\]
Indeed, a finite list of functions $I\to A$ has finitely many
joint value patterns because $A$ is finite; their nonempty
pattern sets form a common partition on which all are constant.
This also proves that $A^I$ is locally finite.

For a finite product of rings, finite projective representations
and nonsingular forms decompose by the central idempotents.
Isometries, orthogonal sums, tensor products and exterior
operations decompose componentwise. Consequently
$\mathcal W(A^\Pi)=\mathcal W(A)^\Pi$.
A refinement map repeats the value on each old cell at all its
new subcells. Apply Theorem~\ref{C:splitWG-filtered-colimits}.
The resulting direct limit is exactly the group of finite-image
functions: every such function is constant on its finite
partition into level sets, and equality of two partition
descriptions is tested on their common refinement.

This also proves injectivity of coordinate evaluation, not
merely its surjectivity. In concrete terms, any projective
idempotent, coaction, Gram matrix and inverse use only finitely
many coefficient functions, so descend to one $A^\Pi$.
Equalities there are precisely coordinatewise equalities.
The asserted canonical description merges all cells with the
same value and omits the zero value. The algorithms intersect
the finitely many level sets, operate on the two finite sets of
canonical values, merge equal output values, and discard empty
or zero cells. Exterior operations require only one set of
values and are treated in the same way. The cases $I=\varnothing$
or $A=0$ give the zero ring on both sides.
\end{proof}

\begin{corollary}[Proved: an explicit non-Noetherian scalar example]
\label{C:finite-ring-powers-scalar}
For $R=\mathbb F_3^I$ and the trivial group $O_0$, every split
quadratic-form class has a unique pair of coordinates $(a,b)$,
where $a:I\to\mathbb Z$ is bounded and
$b:I\to\mathbb F_2$ is arbitrary. Addition is componentwise and
\[
 (a,b)(c,d)=(ac,\overline a\,d+\overline c\,b),\qquad
 \lambda^r(a,b)=
 \left(\binom ar,\,
       b\,\overline{\binom{a-1}{r-1}}\right)\quad(r\geq1).    \tag{C.3.L3}
\]
The bars mean reduction modulo $2$, and all functions and
binomial coefficients are evaluated pointwise; the usual
integer-valued binomial coefficients are used also for negative
integers. For infinite $I$, this is a complete answer over a
non-Noetherian coefficient ring.
\end{corollary}
\begin{proof}
Over $\mathbb F_3$ put
$u=[\langle2\rangle]$ and $t=u-1$.
The finite-field form classification in
Theorem~\ref{C:Om-finite-field-algorithm} gives the additive
group $\mathbb Z\oplus(\mathbb Z/2)t$.
It can also be seen directly by diagonalizing forms: the only
one-dimensional classes are $1,u$, and
$\langle2,2\rangle\simeq\langle1,1\rangle$, while a determinant
comparison proves $u\ne1$ even after stabilization.
Now $u^2=1$, hence $t^2=-2t=0$. Thus
\[
 WG_{\mathrm{sp}}(\mathbb F_3)
       =\mathbb Z[t]/(2t,t^2).
\]
For $b=0$ or $1$ and any $a\in\mathbb Z$,
\[
 \lambda_z(a+bt)
  =(1+z)^a\left(\frac{1+(1+t)z}{1+z}\right)^b
  =(1+z)^a+b\,t z(1+z)^{a-1}.
\]
Taking coefficients gives (C.3.L3). A function into
$\mathbb Z\oplus\mathbb Z/2$ has finite image exactly when
its integer coordinate has finite image, equivalently when it
is bounded. Theorem~\ref{C:finite-ring-powers} proves the result.
For completeness, if $I$ is infinite, choose distinct points
$i_1,i_2,\ldots$. The ideals generated by the first $n$
coordinate characteristic functions form a strictly increasing
chain in $\mathbb F_3^I$, so the ring is not Noetherian.
\end{proof}

\begin{theorem}[Proved: a uniform decision obstruction over locally finite fields]
\label{C:locally-finite-equality-obstruction}
There is no uniform ordinary algorithm which, from a computable
presentation of a locally finite field of odd characteristic
with exact arithmetic and equality, decides all equalities in
$WG_{\mathrm{sp}}(BO_m/K)$, for any fixed $m\geq0$.
This failure already occurs for the two trivial-action forms
$\langle1\rangle$ and $\langle2\rangle$ over uniformly computable
locally finite fields of characteristic $3$.

More precisely, a computable increasing tower of finite fields
can be constructed uniformly from a Turing machine $e$, such
that the difference
$\delta_e=[\langle2\rangle]-[\langle1\rangle]$ is zero in the
union if and only if $e$ halts on its fixed input. The claim is
a uniform obstruction across presentations. It does not assert
that the equality problem is undecidable for every one of these
individual fields. A root-decision or splitting oracle supplies
additional information and is not ruled out by this statement.
\end{theorem}
\begin{proof}
For $s\geq0$ simulate $e$ for $s$ steps and set
\[
 d_{e,s}=
 \begin{cases}
  3^{s+1},&e\text{ has not halted within }s\text{ steps},\\
  2\cdot3^{s+1},&e\text{ has halted within }s\text{ steps}.
 \end{cases}                                             \tag{C.3.L4}
\]
Let $K_{e,s}=\mathbb F_{3^{d_{e,s}}}$.
The degree $d_{e,s}$ divides $d_{e,s+1}$: the quotient is $3$,
except that it is $6$ at the first detected halt.
Use a fixed finite search to construct each finite field, for
example from the first irreducible polynomial of the required
degree, and a fixed finite search for an embedding
$K_{e,s}\hookrightarrow K_{e,s+1}$. These searches terminate
by the existence and subfield properties of finite fields.
No decision about a future halt is used at any stage.
Put $K_e=\mathop{\rm colim}_sK_{e,s}$.
It is an infinite algebraic extension of $\mathbb F_3$ and is
locally finite.

We spell out why this gives computable fields uniformly, not
just enumerable collections of equations. An element may first
be named by a pair consisting of a stage and an element of its
finite field. Equality of two names is decided by embedding
both into their common larger stage. Addition, multiplication
and inverses are computed there as well. For a presentation
with no duplicate names, enumerate stage-element pairs in a
fixed order and retain a pair exactly when it is unequal to
all the finitely many retained earlier pairs. Equality of
names is already decidable, so this is an effective enumeration
of distinct elements. The fields in the tower have unbounded
cardinality, so it produces a domain indexed by all natural
numbers. Compute an operation on this domain using stage
names, then search the enumeration for the name of its result.
That search terminates. All these procedures are uniform in $e$.

In $\mathbb F_{3^d}$ the element $2=-1$ is a square exactly
when $d$ is even. Indeed, the roots of $z^2+1$ lie in
$\mathbb F_9$ and are interchanged by the $3$-power Frobenius;
such a root is fixed by its $d$th iterate precisely for even
$d$. If $e$ halts, an even-degree stage therefore contains a
square root of $2$, giving an isometry
$\langle2\rangle\simeq\langle1\rangle$ and $\delta_e=0$.
If $e$ never halts, all stages have odd degree, and their union
contains no square root of $2$.

Over any field of characteristic different from $2$, the
equality $[\langle2\rangle]=[\langle1\rangle]$ in the split
group would give a stable isometry
$\langle2\rangle\perp T\simeq\langle1\rangle\perp T$ for some
nonsingular form $T$. Determinants show that $2$ must then be a
square. Thus in the nonhalting case $\delta_e\ne0$.
This proves the asserted equivalence for the trivial group.
For any $O_m$, giving a form the trivial action has the
forgetful functor as a retraction; the same equivalence holds
for these trivial-action classes in the equivariant split group.

A uniform equality algorithm in the stated access model,
applied to this computably constructed field and these two
forms, would decide the halting problem. The impossibility
is the usual diagonal undecidability argument, in the
event-detection form of \cite[\S8, p.248]{C-Turing1936}.
The finite-stage criterion of
Theorem~\ref{C:splitWG-filtered-colimits} is fully consistent
with this obstruction: when equality holds, a later finite
witness exists, but no uniform bound or negative decision
follows from that fact.
\end{proof}
